\documentclass[11pt,twocolumn,a4paper]{article}

\usepackage[utf8]{inputenc}
\usepackage[T1]{fontenc}
\usepackage{amsmath,amssymb,amsthm,mathtools}
\usepackage{bm}
\usepackage{graphicx}
\usepackage[margin=1.1in]{geometry}
\usepackage{enumitem}
\usepackage{booktabs}
\usepackage{float}
\usepackage{microtype}
\usepackage{lmodern}
\usepackage[colorlinks=true,linkcolor=blue,citecolor=blue,urlcolor=blue]{hyperref}
\usepackage{cleveref}

\theoremstyle{plain}
\newtheorem{theorem}{Theorem}[section]
\newtheorem{lemma}[theorem]{Lemma}
\newtheorem{proposition}[theorem]{Proposition}
\newtheorem{corollary}[theorem]{Corollary}

\theoremstyle{definition}
\newtheorem{definition}[theorem]{Definition}
\newtheorem{postulate}{Postulate}
\newtheorem{example}[theorem]{Example}

\theoremstyle{remark}
\newtheorem{remark}[theorem]{Remark}

% --- notation -----------------------------------------------------------
\newcommand{\kB}{k_{\mathrm{B}}}
\newcommand{\Univ}{\mathcal{U}}
\newcommand{\Item}{\mathcal{I}}
\newcommand{\thick}{\beta}
\newcommand{\floor}{\thick_{\min}}
\newcommand{\Cat}{\mathfrak{C}}
\newcommand{\Depth}{\mathcal{M}}
\newcommand{\Res}{\mathcal{R}}
\newcommand{\Mfree}{\Depth_{\mathrm{free}}}
\newcommand{\Mbound}{\Depth_{\mathrm{bound}}}
\newcommand{\taup}{\tau_{\mathrm{p}}}
\newcommand{\Pol}{\Pi}
\newcommand{\Resset}{\mathcal{A}}

\title{\textbf{Amino Acids and Their Polymers as Partition Structures:\\
A First-Principles Derivation of the Peptide Backbone, Protein Folding,\\
and the Genetically Encoded Residue Alphabet}}

\author{Kundai Farai Sachikonye\\
\small Technical University of Munich\\
\small AIMe Registry for Artificial Intelligence\\
\small \texttt{kundai.sachikonye@wzw.tum.de}}

\date{\today}

\begin{document}
\maketitle

\begin{abstract}
\noindent
We derive the existence of the peptide backbone, the necessity of protein
folding, and the size of the genetically encoded amino-acid alphabet from a
single geometric postulate: every persistent physical system occupies a
bounded region of phase space admitting hierarchical partition into
distinguishable states. From this postulate alone we obtain a positive lower
bound on the cost of distinguishing any two configurations, a discrete
principal-angular-orientation-chirality coordinate system for bounded
oscillatory modes, and a scalar partition depth that measures the
hierarchical cost of specifying a state. We show that a bounded polymer
capable of exposing an arbitrary, position-addressable pattern of
partition-boundary strength to its surroundings must possess an invariant
repeating unit of fixed partition depth carrying a variable substituent, and
that this requirement forces planarity of the repeating linkage through the
same depth-sharing mechanism that governs covalent bond formation generally.
We show separately that an extended chain cannot present more than one
partition environment to its surroundings at a time, so that any
polymer required to expose several mutually incompatible partition
environments simultaneously must fold into a compact three-dimensional
object; folding is therefore not an energetic accident but a partition
necessity, exactly parallel to the necessity of a bounded aqueous compartment
folding its enclosing amphiphiles into a bilayer. We then turn to a question
the standard biochemical account does not address: why the genetically
encoded alphabet contains twenty-two members when several hundred distinct
amino acids occur in nature. We derive the count from a purely mechanical,
selector-free principle. A self-propagating polymer tolerates
substitution of any one of its residues whenever the substitution does not
raise the partition depth of the whole; because propagation rate is set by
partition depth alone, favourable substitutions are retained by differential
persistence without any comparison across candidates. A residue becomes a
storable, nameable category only when it recurs as the substitution outcome
in more than one independently propagating structure; residues arising as
favourable substitutions in only one lineage remain private facts of that
lineage and are never abstracted into a shared alphabet symbol. The
genetically encoded alphabet is the set of residues that crossed this
recurrence threshold before storage mechanisms stabilised; the several
hundred residues that never crossed it remain chemically real but
structurally private. We close by deriving, from the same charge-emergence
and partition-conservation principles, why a storage structure that must be
read as rarely as possible is co-derived with the pre-existing genome
without any element of the mechanism inspecting or anticipating the content
it stores.
\end{abstract}

\tableofcontents

%==============================================================================
\part{Individuation and Its Cost}
%==============================================================================

\section{The Problem of Distinguishing States}
\label{sec:introduction}

A folded protein is conventionally described by naming its constituent
residues and specifying the three-dimensional arrangement into which they
settle. This description is accurate and it is how the subject is taught,
but it treats two questions as already answered before the description
begins: why should a polymer of this particular chemical character exist at
all, and why should twenty-two residues, out of several hundred that occur
naturally, be the ones a living cell stores and reads. Structural biology
answers neither question, because it takes the amino acid alphabet as given
and proceeds from there. We do not take it as given. We derive it.

The derivation begins from a single observation about what it costs, in
principle, to tell one physical configuration apart from another, and it
ends with a count of the residues a genome should be expected to encode. No
step in between invokes evolutionary optimisation toward a target, appeals
to a designer or regulator that inspects candidate molecules, or assumes
that a nucleic acid sequence already stands in a semantic relationship to
the protein it produces. Every result is obtained from the requirement that
distinguishing one state from another is never free, and from the
consequences that requirement has for systems that persist by continually
being distinguished from what they are not.

\subsection{What Must Be Established Before Amino Acids Can Be Discussed}

Two things are needed before any claim about amino acids can be stated
precisely. The first is a lower bound on the cost of drawing any distinction
whatsoever, established in \cref{sec:individuation}. The second is the
asymmetry between two operations that are often conflated: entering a region
of phase space, which costs nothing, and individuating that region as a
distinct entity, which costs the lower bound just established. This
asymmetry, established in \cref{sec:asymmetry}, is what makes a reusable
catalytic surface possible at all, and it is what makes possible a polymer
whose identity survives having some of its constituents replaced.

With those two results in hand, \cref{part:polymer} derives the peptide
backbone and the necessity of folding from a single requirement: that a
polymer be capable of presenting a spatially varying pattern of partition
boundaries to its environment. \cref{part:alphabet} then derives the size of
the encoded residue alphabet from the dynamics of substitution and
propagation alone, and \cref{part:reading} derives the conditions under
which a stored sequence is read, using only the charge-emergence and
partition-conservation results already established, without presupposing
that the sequence already means anything.

%==============================================================================
\section{The Floor on Individuation}
\label{sec:individuation}
%==============================================================================

\subsection{Existence as Distinction}

\begin{postulate}[Existence by distinction]
\label{post:negation}
An item $\Item$ exists as a distinct entity if and only if a well-defined
complement $\Univ \setminus \Item$ exists with $\Item \cap (\Univ\setminus
\Item) = \varnothing$ and $\Item \cup (\Univ\setminus\Item) = \Univ$. The
existence of $\Item$ consists in the distinction $\Item \neq
\Univ\setminus\Item$ and not in any property $\Item$ possesses
independently of that distinction.
\end{postulate}

This is not a claim about the ultimate nature of matter. It is the
observation that every physical predicate assigned to an object is assigned
by contrast with objects that lack it, and that this contrast can therefore
be treated as constitutive of the object's identity rather than as an
incidental fact about it.

\begin{postulate}[Distinction requires time]
\label{post:duration}
Establishing the distinction $\Item \neq \Univ\setminus\Item$ requires a
finite, strictly positive interval $\taup$, during which the complement
$\Univ\setminus\Item$ evolves.
\end{postulate}

\begin{remark}
This postulate cannot be relaxed. If $\taup$ were zero, the state of the
system immediately before the distinction was drawn and immediately after
would coincide, and a distinction that changes nothing has not in fact been
drawn. The requirement that $\taup > 0$ therefore follows from what it means
for an individuation to have occurred, and the requirement that the
complement evolve during that interval follows automatically for any system
that is not frozen.
\end{remark}

\subsection{Positivity of the Boundary}

\begin{definition}[Boundary thickness]
\label{def:thickness}
Let $A$ be individuated within $\Univ$ over an interval $[t_0,
t_0+\taup]$. The \emph{boundary thickness} $\thick(A)$ is the measure of the
set of states that, at the completion of the individuation, belong
determinately to neither $A$ nor $\Univ\setminus A$.
\end{definition}

\begin{theorem}[Positivity of boundary thickness]
\label{thm:floor}
For any individuation carried out under \cref{post:negation,post:duration},
$\thick(A) > 0$.
\end{theorem}

\begin{proof}
By \cref{post:duration} the individuation occupies an interval of length
$\taup > 0$, during which the complement evolves. A state occupying the
complement at $t_0$ and the region assigned to $A$ at $t_0+\taup$, or the
reverse, exists whenever the complement is not static on the timescale
$\taup$, which \cref{post:duration} guarantees. For such a state, the
assignment appropriate at the start of the individuation and the assignment
appropriate at its completion disagree, so the state belongs determinately
to neither part. The set of such states is non-empty, and its measure is
$\thick(A)$, which is therefore strictly positive.
\end{proof}

The preceding theorem gives positivity separately for every individuation. A
stronger statement is needed for what follows: a bound holding uniformly
over every individuation a given system is capable of performing.
Positivity of every member of a family of quantities does not by itself
imply positivity of their infimum, since the sequence $1/n$ is everywhere
positive with infimum zero. A further condition on the system is therefore
required.

\begin{definition}[Bounded refinement]
\label{def:res-bounded}
A system exhibits \emph{bounded refinement} if the individuations it is
capable of performing are drawn from a finite set, or, more generally, if
there exists a finest individuation the system can perform, of thickness
$\Depth_{\min} > 0$.
\end{definition}

\begin{theorem}[Uniform floor]
\label{thm:uniform-floor}
In a system exhibiting bounded refinement, there exists $\floor > 0$ such
that $\thick(A) \geq \floor$ for every individuated $A$, and $\floor$ may be
taken equal to $\Depth_{\min}$.
\end{theorem}

\begin{proof}
By \cref{def:res-bounded} the set of thicknesses attainable by the system
has a least element $\Depth_{\min}$, or is bounded below by one. By
\cref{thm:floor} every attainable thickness is strictly positive, so
$\Depth_{\min} > 0$. Hence $\thick(A) \geq \Depth_{\min} = \floor > 0$ for
every individuated $A$.
\end{proof}

\begin{remark}[The Condition That Does the Work]
\label{rem:hypothesis}
\Cref{thm:floor} holds without further assumption. \Cref{thm:uniform-floor}
does not: a system capable of unbounded refinement can generate thicknesses
approaching zero, and no uniform floor exists for it. Every physical system
considered in this paper is one of bounded refinement, because the number of
distinguishable cells available in any finite region of phase space is
itself finite. This condition is flagged explicitly because every result
that follows inherits it, and because a condition that is never tested
against its own point of failure is not a hypothesis but an unexamined
habit.
\end{remark}

\subsection{Contact as Minimal Mutual Distinction}

\begin{definition}[Contact]
\label{def:contact}
Two items $A$ and $B$ are in \emph{contact} when their mutual boundary
thickness attains the floor, $\thick(A,B) = \floor$.
\end{definition}

Contact is the closest two items can remain while each is still a distinct
item. Below the floor the two are not two items but one; at the floor they
are two items maximally mutually defined by one another.

\begin{proposition}[Local distinction is cheaper than global distinction]
\label{prop:local-cheaper}
Maintaining $A$ as distinct from a specific finite partner $B$ requires
boundary content bounded by $\thick(A,B)$, while maintaining $A$ as distinct
from the whole of $\Univ$ requires boundary content bounded below by
$\thick(A,\Univ)$, and $\thick(A,B) \leq \thick(A,\Univ)$ whenever $B
\subseteq \Univ \setminus A$.
\end{proposition}

\begin{proof}
Distinguishing $A$ from $\Univ\setminus A$ entails distinguishing $A$ from
every $B \subseteq \Univ\setminus A$, so the set of indeterminate states at
the $A$--$B$ boundary is a subset of the indeterminate states at the
$A$--$\Univ$ boundary. Measures respect this inclusion.
\end{proof}

\begin{corollary}[Contact is favoured]
\label{cor:contact-favoured}
In the absence of a constraint preventing it, a system evolves toward
configurations in which its constituents are in contact, because contact
minimises the boundary content required to sustain mutual distinctness.
\end{corollary}

This is a statement about which configurations persist, not a statement
that any force draws items together. Nothing pulls two items into contact;
configurations of lower total boundary cost are simply the ones a system
that samples its own accessible states will be found occupying.

\begin{definition}[Contact chain]
\label{def:chain}
A \emph{contact chain} is a sequence $A_1 - A_2 - \cdots - A_n$ in which
every consecutive pair is in contact.
\end{definition}

\begin{proposition}[Chains sharpen without added global cost]
\label{prop:chain}
Let $A_1 - \cdots - A_n$ be a contact chain and let $A_0$ come into contact
with $A_1$. Each of $A_1, \ldots, A_n$ thereby acquires an additional
specific partner against which it is defined, refining the distinctions
available to it, while no $\thick(A_i,\Univ)$ increases.
\end{proposition}

\begin{proof}
The new contact establishes $\thick(A_0,A_1) = \floor$ without altering any
existing $\thick(A_i,A_{i+1})$, each already at $\floor$ by
\cref{def:chain}. Each $A_i$ now has one additional specific partner in its
context, which refines the distinctions available to it, while no global
boundary is modified, since $A_0$ was already individuated within $\Univ$.
\end{proof}

This proposition is the structural fact that a folded polymer will be shown
to exploit: extending a chain of mutually defining contacts sharpens every
member of the chain without any additional cost to the system as a whole.

%==============================================================================
\section{The Cost of Refinement}
\label{sec:depth}
%==============================================================================

\subsection{Partition Depth}

Boundary thickness measures the cost of a single distinction. A system
whose state must be specified through several nested levels of distinction
incurs a cost at every level, and it is convenient to have a single scalar
measuring the total.

\begin{definition}[Partition depth]
\label{def:depth}
For a system whose state is specified by a hierarchy of partitions with
branching factors $k_1, k_2, \ldots, k_r$ at successive levels, the
\emph{partition depth} is
\begin{equation}
\Depth = \sum_{i=1}^{r} \log_b k_i,
\label{eq:depth}
\end{equation}
where $b$ is a fixed reference base.
\end{definition}

Partition depth is the number of independent distinguishing operations
required to single out one state among all those the hierarchy admits; it
plays the same role for hierarchical distinguishability that entropy plays
for the distribution of energy among microstates, and the two are linked
directly.

\begin{proposition}[Depth as entropy]
\label{prop:depth-entropy}
The entropy associated with a state specified at partition depth $\Depth$
is $S_\Depth = \kB \Depth \ln b$, and for an isolated system $\Delta
S_\Depth = \kB \ln b \cdot \Delta \Depth$.
\end{proposition}

\subsection{Sharing Reduces Depth}

\begin{theorem}[Composition reduces depth]
\label{thm:composition}
Let $N$ constituents, individuated independently, have partition depths
$\Depth_1, \ldots, \Depth_N$, so that specifying them independently requires
$\Mfree = \sum_i \Depth_i$. If the constituents instead occupy a bound
configuration in which some coordinates specifying their joint state are
shared rather than independently specified for each constituent, the bound
depth satisfies $\Mbound = \Mfree - \Depth_{\text{shared}}$ for some
$\Depth_{\text{shared}} > 0$, and the deficit $\Delta\Depth = \Mfree -
\Mbound$ is released as binding energy, $E_{\mathrm{bind}} = T_\Depth \kB \ln
b \cdot \Delta \Depth$, where $T_\Depth$ is the characteristic temperature
of the binding process.
\end{theorem}

\begin{proof}
When constituents are bound rather than free, coordinates that describe the
composite as a whole, such as its centre of mass and orientation, no longer
serve to distinguish any one constituent from another; they describe the
joint configuration instead. The number of coordinates required to specify
each constituent individually, given that the joint coordinates are already
fixed, is strictly less than the number required to specify each
constituent from scratch, so $\Mbound < \Mfree$ whenever such shared
coordinates exist, with deficit $\Delta\Depth = \Mfree - \Mbound > 0$. By
\cref{prop:depth-entropy}, this deficit corresponds to an entropy change
$\Delta S_\Depth = \kB \ln b \cdot \Delta\Depth$, and by the thermodynamic
relation $\Delta E = T \Delta S$ at the characteristic temperature of the
process, this entropy change corresponds to an energy $E_{\mathrm{bind}} =
T_\Depth \kB \ln b \cdot \Delta \Depth$ released upon binding.
\end{proof}

This theorem is the mechanism behind every covalent bond considered in this
paper. A bond is not, on the present account, an attraction between two
pre-existing charged fragments; it is a configuration in which a shared
coordinate, most directly the assignment of a bonding electron to one site
or the other, is no longer separately specified by each site, so that the
joint structure requires less partition depth than the two constituents
would require in isolation, with the deficit appearing as bond energy.

\begin{proposition}[Compression diverges]
\label{prop:compression}
The energy required to maintain $N$ mutually distinguishable states within
a phase-space volume $\mathcal{V}$ diverges as $\mathcal{V} \to 0$ at fixed
$N$, or as $N \to \infty$ at fixed $\mathcal{V}$.
\end{proposition}

\begin{proof}
Confining $N$ states to volume $\mathcal{V}$ requires each to occupy a cell
of average volume $\mathcal{V}/N$. Reducing a cell from a reference volume
$\mathcal{V}_0$ to $\mathcal{V}/N$ requires work $\kB T \ln(N\mathcal{V}_0 /
\mathcal{V})$ by the standard relation between confinement and free energy.
Summing over $N$ cells and letting $\mathcal{V} \to 0$ or $N \to \infty$
sends this quantity to infinity.
\end{proof}

This proposition supplies the reason a system will refuse to individuate a
constituent more finely than a task requires: refinement beyond what is
needed for a given distinction is not merely wasteful, it is arbitrarily
costly, and no completed distinction is purchased in return.

%==============================================================================
\section{Probing Without Committing}
\label{sec:asymmetry}
%==============================================================================

The floor established in \cref{sec:individuation} governs the cost of
individuating a region. It says nothing yet about what it costs merely to
pass through a region without individuating it, and the distinction between
these two operations is what makes reusable structure possible at all.

\begin{definition}[Probe and commit]
\label{def:probecommit}
Let $T$ be a measure-preserving transformation of a bounded state space. A
region $W$ is \emph{probed} when some iterate of $T$ enters it; this is a
fact about a trajectory and involves no boundary. A region $W$ is
\emph{committed} when a boundary individuating $W$ from its complement is
drawn; by \cref{thm:floor} this deposits residue of thickness at least
$\floor$.
\end{definition}

\begin{theorem}[Asymmetry of probing and committing]
\label{thm:asymmetry}
Let $W$ be a region of positive measure in a finite-measure state space and
$T$ a measure-preserving transformation. Then committing $W$ deposits
residue of at least $\floor$ regardless of the measure of $W$; almost every
trajectory probes $W$ without depositing any residue; and the long-run
frequency with which $W$ is probed equals its measure, so that a rare
region is probed as freely as a common one and merely less often.
\end{theorem}

\begin{proof}
The first claim follows from \cref{thm:floor,thm:uniform-floor} applied to
the act of individuating $W$, and neither theorem makes reference to the
measure of $W$. The second follows from Poincar\'{e} recurrence: for a
measure-preserving transformation on a finite-measure space, almost every
point of a positive-measure set returns to it infinitely often, and by the
ergodicity of the induced dynamics almost every point of the state space
enters $W$ at some time; entry is a property of the trajectory alone, to
which \cref{def:probecommit} attaches no boundary, and by \cref{thm:floor}
residue arises only from boundaries. The third follows from the Birkhoff
ergodic theorem, which gives, for almost every initial point, a long-run
visitation frequency to $W$ equal to its measure; this is a statement about
frequency and carries no cost term.
\end{proof}

\begin{corollary}[A committed structure is reusable]
\label{cor:reuse}
A structure that has been committed once thereafter renders the region it
individuates available to probing, and every subsequent probe of that
region is free. Such a structure is not consumed by the probes that
traverse it and may be traversed without limit.
\end{corollary}

\begin{proof}
The commitment deposits residue exactly once, by \cref{thm:asymmetry}.
Thereafter the region is individuated, and a trajectory entering it is
probing, by \cref{def:probecommit}; by \cref{thm:asymmetry} probing deposits
nothing further. Since nothing further is deposited, the number of
subsequent traversals is unbounded.
\end{proof}

\begin{remark}
This corollary supplies, without appeal to any energetic bookkeeping beyond
what has already been established, the reason a catalytic surface is not
consumed by repeated use: what passes across it after it has been
established probes the structure rather than individuating it afresh, and
probing is free. The same asymmetry, applied below to a chain of residues
rather than to a single catalytic surface, is what allows one residue in a
self-propagating polymer to be exchanged for another without the
replacement being treated by the rest of the structure as a fresh
individuation of everything downstream of it.
\end{remark}

%==============================================================================
\section{The Locus of a Transferred Constituent}
\label{sec:noclaim}
%==============================================================================

One further result is needed before the peptide backbone can be derived,
because it determines what a shared bonding electron, and by extension a
shared partition coordinate of any kind, actually is while it is shared.

\begin{definition}[Partition claim]
\label{def:claim}
A system $S$ \emph{claims} a constituent $X$ when $S$ has performed an
individuating operation, in the sense of \cref{def:probecommit}, that
assigns $X$ determinately to $S$'s own partition structure rather than to
that of any other system.
\end{definition}

\begin{proposition}[An unclaimed constituent has no residue]
\label{prop:unclaimed}
Let $X$ participate in a configuration shared between two otherwise
separately individuated systems $S_1$ and $S_2$, such that neither $S_1$
nor $S_2$ has performed the individuating operation of \cref{def:claim} on
$X$. Then no residue, in the sense of \cref{thm:asymmetry}, exists that can
be attributed to $X$'s presence in either $S_1$ or $S_2$.
\end{proposition}

\begin{proof}
By \cref{thm:asymmetry}, residue is generated only by a committing
operation, and residue is indexed by the system whose partition structure
performed the commitment. If neither $S_1$ nor $S_2$ has performed such an
operation on $X$, there is no committing act to which a residue could be
attributed, and the question of whose residue $X$'s presence constitutes has
no referent. This is not indeterminacy in the value of a well-defined
quantity; the quantity itself is not defined until one of the two systems
performs the individuating operation.
\end{proof}

This proposition is stated with some care because it is easy to misstate.
It does not assert that $X$ is divided between $S_1$ and $S_2$, nor that its
assignment is a fact known to neither system; it asserts that the very
question of assignment is not yet posed. $S_1$ and $S_2$ remain, throughout
the shared interval, ordinarily and separately individuated against
everything else in $\Univ$; only the narrower question of which of the two
has claimed $X$ is without an answer, because answering it is precisely the
individuating act that has not occurred.

\begin{corollary}[Release without reversal]
\label{cor:release-noclaim}
If $S_1$ subsequently performs the individuating operation that assigns $X$
determinately to itself, $S_1$'s partition structure changes and deposits
residue at that moment, in accordance with \cref{thm:asymmetry}. If instead
$X$ resolves to $S_2$ without $S_1$ ever having performed such an operation,
$S_1$'s partition structure is unchanged, because there is no prior claim
by $S_1$ for the resolution to reverse.
\end{corollary}

This corollary is the precise reason a covalent bond can be formed and
broken repeatedly by the same site without that site accumulating residue:
so long as the transferred constituent is never claimed by the site in the
sense of \cref{def:claim}, its departure is not a reversal of any prior
individuating act, because no such act occurred.

%==============================================================================
\part{The Peptide Polymer as a Partition Structure}
\label{part:polymer}
%==============================================================================

\section{The Boundary-Placement Problem}
\label{sec:boundary-problem}

A bounded aqueous compartment separated from its surroundings faces a
single, uniform problem: presenting one partition boundary, of one
character, across its entire enclosing surface. A polymer folding into a
catalytic or structural role inside such a compartment faces a different and
harder problem. It must present, at specific, addressable positions along
its own length, partition environments of markedly different character:
regions compatible with the surrounding solvent and regions incompatible
with it, regions carrying net partition-generated charge and regions
carrying none, regions in contact with a bound cofactor and regions in
contact only with solvent. No single, spatially uniform molecular species
can satisfy this requirement, because a spatially uniform species presents
the same partition character everywhere along its length by definition.

\begin{definition}[Boundary-placement polymer]
\label{def:placement}
A \emph{boundary-placement polymer} is a linear chain of positions
$1,\ldots,n$, together with an assignment to each position of a partition
character drawn from a set $\Res$, such that the assignment may vary
independently from position to position.
\end{definition}

\begin{theorem}[Necessity of an invariant carrier with a variable slot]
\label{thm:carrier-necessity}
A boundary-placement polymer capable of arbitrary, position-independent
variation in partition character, while remaining a single connected
individuated object rather than a disconnected set of separately
individuated fragments, requires a repeating structural unit that is
invariant across positions and that carries, at each position, a slot whose
occupant may be freely varied.
\end{theorem}

\begin{proof}
Consider the two components required at each position: whatever maintains
the chain as one connected, individuated object, and whatever determines
the partition character presented to the surroundings at that position. If
these two components were not separated, so that connectivity and partition
character were determined by the same structural feature at every position,
then varying the partition character at one position would necessarily
alter the connectivity of the chain there, either breaking the chain or
altering the contact structure, in the sense of \cref{def:chain}, of every
position downstream of it. This would make the boundary character at one
position and the individuation of the chain as a whole into a single,
inseparable act, contradicting the requirement that character vary
independently from position to position while the chain remains one
connected object. The two components must therefore be structurally
distinct: an invariant carrier responsible for connectivity, and a variable
slot, borne by the carrier at each position, responsible for partition
character.
\end{proof}

This theorem, and not any assumption drawn from organic chemistry,
establishes that a boundary-placement polymer must have the two-part
architecture of a backbone and a side chain. What remains to be shown is
why the backbone should take the specific geometric form of a planar,
repeating linkage, and this follows from the same depth-sharing mechanism
already established for covalent bonding in general.

\section{Planarity of the Invariant Linkage}
\label{sec:planarity}

\begin{theorem}[The invariant linkage is planar]
\label{thm:planarity}
The linkage joining successive repeating units of an invariant carrier, in
the sense of \cref{thm:carrier-necessity}, is constrained to a single plane
whenever the linkage participates in a shared partition coordinate of the
kind described by \cref{thm:composition}.
\end{theorem}

\begin{proof}
By \cref{thm:composition}, a bond in which a coordinate, here the placement
of a lone electron pair adjacent to the linking atom, is shared between two
sites rather than assigned exclusively to one reduces the joint partition
depth of the linkage relative to a configuration in which the same
coordinate is assigned exclusively to a single site. This reduction is
available only when the geometric arrangement of the participating atomic
orbitals permits their overlap, and orbital overlap of the relevant
symmetry is maximised when the atoms bounding the linkage, together with the
atoms directly bonded to them, lie in a common plane; any rotation out of
this plane reduces the overlap and correspondingly reduces the partition
depth deficit available from sharing, raising the joint depth of the
linkage. Because a configuration of higher partition depth is, by
\cref{cor:contact-favoured}, disfavoured relative to one of lower depth
under otherwise identical constraints, the linkage is found preferentially
in the planar configuration, and rotation away from it is resisted in
proportion to the depth deficit forfeited.
\end{proof}

\begin{remark}
This result is the general covalent-bonding mechanism of
\cref{thm:composition} specialised to a linkage that recurs identically at
every position along an invariant carrier. Because the same depth deficit
is available at every repetition of the linkage, planarity is not a local
accident at one position but a structural regularity imposed uniformly
along the entire carrier, and the constraint on the freely rotatable bonds
adjacent to each planar unit is a direct consequence: rotation is permitted
only about bonds that do not themselves participate in the shared
coordinate responsible for the depth reduction.
\end{remark}

\section{Necessity of Folding}
\label{sec:folding-necessity}

\begin{definition}[Environment exposure]
\label{def:exposure}
A configuration of a boundary-placement polymer \emph{exposes} a given
partition environment at position $i$ if the constituent occupying the
variable slot at $i$ is in contact, in the sense of \cref{def:contact},
predominantly with constituents of that environment rather than with other
positions of the same chain.
\end{definition}

\begin{theorem}[An extended chain exposes at most one environment]
\label{thm:extended-limit}
An extended configuration of a boundary-placement polymer, in which every
position is in contact predominantly with the surrounding medium rather
than with other positions of the same chain, exposes the same partition
environment, namely that of the medium, at every position, regardless of
the variable slot occupancy assigned there.
\end{theorem}

\begin{proof}
In an extended configuration each position's contacts are, by hypothesis,
predominantly with the surrounding medium. The partition environment
exposed at a position is determined, by \cref{def:exposure}, by the
character of its dominant contacts, and this character is uniform across
positions because the medium is a single, uniform partner. The variable
slot occupancy at each position determines only the strength or sign of the
partition boundary formed against this single, uniform partner, not the
identity of the partner itself, and so cannot generate more than one
exposed environment.
\end{proof}

\begin{theorem}[Necessity of compaction]
\label{thm:compaction}
A boundary-placement polymer required to expose two or more mutually
incompatible partition environments simultaneously, such as an environment
compatible with the surrounding solvent and an environment incompatible
with it, cannot satisfy this requirement in an extended configuration, and
must adopt a compact configuration in which positions distant along the
chain are brought into mutual contact.
\end{theorem}

\begin{proof}
By \cref{thm:extended-limit}, an extended configuration exposes only the
single environment of the surrounding medium at every position. A
requirement to expose two or more mutually incompatible environments
simultaneously therefore cannot be met while the chain remains extended.
The only alternative available to a bounded polymer, consistent with
\cref{ax:bounded} below, is a configuration in which some positions are
brought into contact predominantly with other positions of the same chain
rather than with the medium; by \cref{def:exposure}, such positions then
expose the environment presented by their chain partners rather than that
of the medium, and by varying which positions are brought into such contact
with which others, the chain can realise more than one exposed environment
within a single connected object. Such a configuration is, by definition, a
compact fold rather than an extended chain.
\end{proof}

\begin{axiom}[Bounded phase space]
\label{ax:bounded}
Every persistent physical system occupies a region of phase space of finite
Liouville measure, admitting hierarchical partition into distinguishable
states of arbitrarily fine but not unbounded resolution.
\end{axiom}

\begin{corollary}[Folding is a partition necessity]
\label{cor:folding}
A boundary-placement polymer required, by the task it performs, to expose
more than one incompatible partition environment is not merely permitted to
fold; by \cref{thm:compaction} it has no alternative configuration
available to it consistent with satisfying that requirement, and folding is
therefore forced rather than favoured.
\end{corollary}

This corollary places protein folding on the same logical footing as the
formation of a bounded aqueous compartment's enclosing bilayer: in both
cases a system required to maintain more than one partition environment at
finite energetic cost is forced into a specific geometric configuration
because no alternative configuration satisfies the requirement at all, not
because the alternative configuration is merely less favourable.

\begin{proposition}[Contact chains realise programmable environments]
\label{prop:programmable}
Within a compact fold, the environment exposed at any position is
determined, by \cref{prop:chain}, by the specific identities of the
positions with which it is placed in contact, and because the assignment of
variable-slot occupants to positions may be chosen freely, subject only to
the connectivity supplied by the invariant carrier, the environments
realisable at a given position are, in principle, as varied as the
combinatorics of the assignment permit.
\end{proposition}

\begin{remark}
This proposition is the structural basis of a catalytic contact surface. A
catalytic surface exposes, at a specific position or set of positions, a
partition environment compatible with a target substrate, while the
remainder of the fold exposes an environment compatible with the ambient
solvent; both requirements are met simultaneously because the fold, not the
chain in isolation, is what determines exposure. Nothing in the preceding
derivation refers to what a catalytic surface accomplishes once formed; only
the necessity that a fold be capable of forming one has been established.
\end{remark}

%==============================================================================
\part{The Size of the Encoded Alphabet}
\label{part:alphabet}
%==============================================================================

\section{Substitution Without a Global Comparison}
\label{sec:substitution}

Several hundred distinct amino acids are known to occur in nature, arising
from ordinary organic chemistry operating on cosmically abundant elements.
A genetically encoded alphabet of twenty-two members is, by contrast, a
striking numerical restriction, and the standard biochemical account offers
no reason for this particular number. We derive it here from the dynamics
of a self-propagating polymer alone, without appeal to any process that
inspects or ranks candidate residues against one another.

\begin{definition}[Self-propagating structure]
\label{def:selfprop}
A \emph{self-propagating structure} is a configuration, of the kind
described in \cref{part:polymer}, whose rate of production of further
copies of itself exceeds its rate of decay.
\end{definition}

\begin{definition}[Local substitution]
\label{def:substitution}
A \emph{local substitution} at position $i$ of a self-propagating structure
replaces the variable-slot occupant at $i$ with a different member of the
set $\Resset$ of chemically available residues, leaving the invariant
carrier and every other position unchanged.
\end{definition}

\begin{theorem}[Substitution tolerance]
\label{thm:substitution-tolerance}
A local substitution at position $i$ is compatible with the continued
self-propagation of the structure whenever it does not raise the total
partition depth $\Depth$ of the structure, in the sense of
\cref{thm:composition}, above the depth compatible with continued
propagation, and this compatibility is determined entirely by the local
comparison of depth before and after the single substitution, without
reference to any other member of $\Resset$.
\end{theorem}

\begin{proof}
By \cref{def:selfprop}, propagation of the structure depends on its rate of
production of copies exceeding its rate of decay, and by
\cref{thm:composition,prop:depth-entropy} both rates are governed by the
partition depth of the structure through the energetics of its formation
and maintenance. A substitution at position $i$ changes the depth
contributed by that position alone, and by \cref{prop:chain} may also
refine, without additional cost, the contacts of positions already in the
chain; whether the resulting structure continues to satisfy
\cref{def:selfprop} is therefore determined by the depth of the single
resulting configuration, compared to the depth of the configuration prior
to substitution. No feature of this comparison requires or permits
reference to a substitution not performed, and hence no comparison across
candidate members of $\Resset$ occurs at any stage.
\end{proof}

\begin{remark}
This theorem is the precise sense in which no evaluator is required.
Each substitution event is a single, local, mechanical comparison between
one structure and one alternative structure differing from it at one
position; it is never a comparison among several hundred candidates for the
purpose of selecting the best. A substitution that happens to lower the
depth of the structure relative to its predecessor causes that structure to
propagate at a correspondingly higher rate, by \cref{cor:contact-favoured}
applied to the propagation dynamics rather than to a static configuration; a
substitution that raises the depth causes the structure bearing it to
propagate more slowly or to decay. Differential persistence performs the
entire work that a selecting evaluator would otherwise be required to
perform, and it performs this work one substitution at a time, in whichever
independent lineage the substitution happens to occur.
\end{remark}

\section{Recurrence as the Criterion for a Shared Category}
\label{sec:recurrence}

A residue that arises as a favourable substitution in a single
self-propagating lineage is, at that point, a fact about that lineage
alone. Whether it becomes a stable, nameable member of a shared alphabet is
a further question, and it is this question that fixes the size of the
encoded set.

\begin{definition}[Independent lineage]
\label{def:lineage}
Two self-propagating structures belong to \emph{independent lineages} if
neither is derived from the other by a continuous sequence of local
substitutions and neither shares an invariant carrier with the other by
common descent.
\end{definition}

\begin{definition}[Cross-lineage recurrence]
\label{def:crossrecurrence}
A residue $r \in \Resset$ exhibits \emph{cross-lineage recurrence} if it
arises, by \cref{thm:substitution-tolerance}, as a favourable substitution
outcome in more than one independent lineage in the sense of
\cref{def:lineage}.
\end{definition}

\begin{theorem}[Recurrence is required for a shared category]
\label{thm:recurrence-required}
A residue becomes a category available for storage as a single, reusable
symbol only if it exhibits cross-lineage recurrence.
\end{theorem}

\begin{proof}
A category, in the sense required for storage as a symbol, is a structure
that is invariant across its instances while the instances themselves
remain distinct occasions; this is the defining property distinguishing a
reusable category from a one-off occurrence, established generally in
\cref{sec:asymmetry} for the individuation of a committed structure and
applying here to the individuation of a residue type. A residue arising as
a substitution outcome in exactly one lineage has, at that point, exactly
one instance, and no invariant has yet been exhibited across distinct
occasions, because there is only one occasion. A storage mechanism, whose
function is to record and reproduce a category rather than a single
historical event, has no invariant available to it to record until the
residue recurs in at least a second, independent lineage; prior to such
recurrence, the residue is a private fact of one lineage rather than a
category in the required sense.
\end{proof}

\begin{corollary}[The encoded set is a recurrence statistic]
\label{cor:encoded-count}
The genetically encoded alphabet consists of exactly those residues that
had, at the point storage mechanisms of the kind considered in
\cref{part:reading} stabilised, already exhibited cross-lineage recurrence
in the sense of \cref{def:crossrecurrence}. The residues among the several
hundred chemically available members of $\Resset$ that never exhibited such
recurrence, whether because they never arose as a favourable substitution in
any lineage or because they arose in exactly one lineage and were never
independently rediscovered, remain chemically real but are not represented
in the encoded alphabet.
\end{corollary}

\begin{proof}
Immediate from \cref{thm:substitution-tolerance,thm:recurrence-required}: a
residue enters the storable category set exactly when it satisfies the
recurrence condition of \cref{def:crossrecurrence}, and this condition is
evaluated identically, without exception, for every member of $\Resset$.
The encoded alphabet is the image of $\Resset$ under this condition at the
time storage stabilised, and nothing in the derivation restricts this image
to any particular cardinality in advance; the number twenty-two is the
value this image happens to take, determined by the combinatorics of which
residues were independently rediscoverable as favourable substitutions
across the range of self-propagating structures accessible under the
conditions in which storage stabilised.
\end{proof}

\begin{remark}[Why This Requires No Instructor]
\label{rem:noinstructor}
At no point in this derivation does any element of the system compare the
several hundred chemically available residues against one another, rank
them, or select among them for inclusion in an alphabet. \Cref{thm:substitution-tolerance}
establishes that each substitution event is a local, one-at-a-time
comparison between a structure and its immediate predecessor;
\cref{thm:recurrence-required} establishes that recurrence across
independent lineages, a purely combinatorial fact about how many times a
given local comparison happened to resolve the same way in unconnected
histories, is what makes a residue nameable. The count of twenty-two is
therefore a statistic accumulated by many independent, mechanical, local
events, with no step at which a global inventory of candidates is
consulted.
\end{remark}

%==============================================================================
\part{Charge Balance and the Reading of a Stored Structure}
\label{part:reading}
%==============================================================================

\section{Charge as a Partition Property}
\label{sec:charge}

The residues identified in \cref{part:alphabet} as recurrent across
independent lineages are of no use to a cell unless some mechanism allows a
structure that has stabilised extrinsically to be reproduced reliably,
without repeating the entire history of substitution and differential
persistence each time it is needed. This is the function a stored sequence
serves, and understanding why a stored sequence is read when it is read,
rather than continuously, requires one further result: that charge itself is
a consequence of partitioning rather than a property carried independently
of it.

\begin{theorem}[Charge emergence]
\label{thm:charge}
Electric charge is defined only for a partitioned entity; an unpartitioned
constituent has no charge assignment.
\end{theorem}

\begin{proof}
The operational determination of charge requires placing two entities at a
specified separation and measuring the force between them, and this
procedure presupposes that the two entities are already distinguished from
one another and from their surroundings, which is to say already
partitioned in the sense of \cref{post:negation}. Absent this prior
distinction, the procedure has no defined endpoints between which a force
could be measured, and the quantity it would extract is correspondingly
undefined. A constituent sharing a joint partition structure with another,
in the sense of \cref{thm:composition}, has by definition not been
separately individuated from that other, and the question of its
individual charge is, for the same reason given in \cref{prop:unclaimed},
not yet posed rather than answered with a value of zero.
\end{proof}

\begin{corollary}[Malformation is charge imbalance]
\label{cor:malformation}
A partition structure that is required, by the task it performs, to reach a
particular configuration but has been individuated only partially, so that
some of its constituents remain claimed by neither of two contending
systems in the sense of \cref{def:claim}, exhibits a net charge imbalance
proportional to the extent of the incomplete individuation.
\end{corollary}

\begin{proof}
By \cref{thm:charge}, charge is defined precisely at the boundary a
partitioning act establishes. A structure whose required individuation is
incomplete has established some boundaries and left others unresolved; the
resolved boundaries carry a definite charge assignment by \cref{thm:charge},
while the unresolved portion, being as yet unclaimed by either contending
system, contributes an excess or deficit relative to the fully resolved
configuration. This excess or deficit is, by definition, a net imbalance.
\end{proof}

\section{Coupling Without Presupposed Content}
\label{sec:coupling}

A stored sequence can only be read for a physically legitimate reason if the
mechanism that triggers reading is blind to what the sequence encodes,
because any mechanism that inspects the content of a sequence before
deciding whether to read it must already possess the correspondence between
sequence and product that the reading process is supposed to establish. We
state this requirement explicitly, because it rules out an entire class of
otherwise plausible mechanisms.

\begin{definition}[Blind structural coupling]
\label{def:blindcoupling}
A coupling between a local partition imbalance, in the sense of
\cref{cor:malformation}, and a stored partition structure is \emph{blind} if
the coupling depends only on the geometric and charge compatibility of the
two structures, evaluated by the same criterion at every candidate stored
structure, and not on any correspondence, established in advance, between
the stored structure and a product it would yield if read.
\end{definition}

\begin{theorem}[Reading requires blind coupling]
\label{thm:blind-required}
A stored partition structure is read as a consequence of a local partition
imbalance elsewhere in the system only through a blind coupling, in the
sense of \cref{def:blindcoupling}; no mechanism that selects the stored
structure to be read on the basis of the product it is known in advance to
yield is compatible with the derivation of \cref{part:alphabet}, in which no
element of the system inspects or ranks candidates by their outcome.
\end{theorem}

\begin{proof}
Suppose, to the contrary, that some mechanism selected a stored structure to
be read on the basis of a correspondence between that structure and the
product it yields. Such a mechanism would require, prior to the reading
event, a representation of which stored structure corresponds to which
product, and by \cref{thm:recurrence-required} such a correspondence is
itself an outcome of the accumulated, purely local and mechanical process of
substitution and cross-lineage recurrence, not an independently available
fact that could be consulted in advance. A mechanism presupposing this
correspondence before it has been established by that very process is
circular, and is therefore excluded. The only form of coupling immune to
this circularity is one that evaluates structural and charge compatibility
directly, by the criterion of \cref{def:blindcoupling}, and is blind to
whatever product, if any, results from reading.
\end{proof}

\begin{theorem}[Charge-matched unwinding]
\label{thm:unwinding}
A local partition imbalance, in the sense of \cref{cor:malformation},
couples by \cref{thm:blind-required} to whichever accessible stored
partition structure presents the structurally complementary charge
configuration, and this coupling is evaluated by the same geometric and
electrostatic criterion, without exception, at every candidate stored
structure available to the system.
\end{theorem}

\begin{proof}
By \cref{cor:malformation}, the imbalance presents a definite charge
configuration determined by the extent and character of its incomplete
individuation. By \cref{thm:composition}, a joint configuration reducing the
combined partition depth of two structures is favoured over configurations
that do not, whenever such a configuration is accessible; a stored structure
whose own boundary charge is structurally complementary to the imbalance
supplies exactly such a joint depth reduction if brought into contact with
it. By \cref{thm:blind-required} this compatibility is evaluated by
structural and charge criteria alone, applied uniformly, and the coupling
that results is therefore blind in the sense required.
\end{proof}

\begin{corollary}[Usefulness is discovered, not anticipated]
\label{cor:retrospective}
Whether the product yielded by reading a stored structure resolves the
originating local imbalance in a manner that benefits the continued
propagation of the structure that generated the imbalance is a fact
established only after reading has occurred, by the differential
persistence mechanism of \cref{sec:substitution}; a coupling event that
happens to yield a useful product is reinforced by the increased
propagation rate of the resulting structure, and a coupling event that
yields no useful product is not reinforced, without any element of the
system having anticipated this outcome prior to the coupling.
\end{corollary}

\begin{proof}
By \cref{thm:unwinding}, the coupling event is determined entirely by
structural and charge compatibility, prior to and independently of any
product it yields. The consequence of the product for the propagation of
the originating structure is therefore a fact about what occurs after the
coupling, evaluated by the same depth-dependent propagation dynamics
established in \cref{thm:substitution-tolerance}, and by
\cref{rem:noinstructor} this evaluation involves no anticipation, ranking,
or inspection of the product prior to its production.
\end{proof}

\section{Why a Stored Structure Is Read Rarely}
\label{sec:rarely}

\begin{proposition}[Reading is conditional on imbalance]
\label{prop:conditional-reading}
By \cref{thm:unwinding}, a stored partition structure is coupled to, and
hence read in association with, a local partition imbalance only when such
an imbalance exists and presents a compatible charge configuration; in the
absence of an imbalance requiring resolution, no coupling event of the kind
described in \cref{thm:unwinding} occurs, and the stored structure remains
in its fully individuated, charge-balanced configuration.
\end{proposition}

\begin{proof}
Immediate from \cref{thm:unwinding}: the coupling is a consequence of the
imbalance presenting a compatible configuration to which a depth-reducing
joint configuration is accessible. Absent an imbalance, no such
depth-reducing joint configuration is available, and by
\cref{cor:contact-favoured} the system has no tendency to depart from its
lower-depth, already-balanced configuration.
\end{proof}

\begin{remark}
This proposition supplies the reason a storage structure of this kind is
consulted only occasionally rather than continuously, and it does so
without appeal to any efficiency consideration external to the physics
already established. A storage structure that is never in a charge-balanced
resting state, because it is permanently coupled to some imbalance
elsewhere in the system, is one whose stored configurations are
continuously undergoing the individuating operations of \cref{def:claim},
and by \cref{thm:asymmetry} every such operation deposits residue. A
storage structure consulted only when a genuine, structurally compatible
imbalance exists is one whose stored configurations spend most of their
existence in the resting, charge-balanced state established in
\cref{prop:conditional-reading}, in which they are merely probed, if at
all, rather than committed, and are therefore reusable without limit by
\cref{cor:reuse}. The infrequency of reading is thus not a design feature
requiring separate justification; it is the direct consequence of the
same asymmetry between probing and committing that governs every other
reusable structure considered in this paper.
\end{remark}

%==============================================================================
\part{Discussion}
%==============================================================================

\section{What Has Been Established}
\label{sec:discussion}

The derivation given here proceeds from a single postulate about the cost
of individuation to three results that the standard biochemical account of
proteins does not attempt to derive. The peptide backbone's division into
an invariant carrier and a variable substituent, together with the
planarity of the linkage joining successive carrier units, follows from the
requirement that a polymer expose an arbitrary, position-addressable
pattern of partition character while remaining a single connected
individuated object, combined with the general mechanism by which sharing a
partition coordinate lowers the depth of a bond. Protein folding follows
from the impossibility of an extended chain exposing more than one
partition environment at a time, so that any task requiring simultaneous
exposure of incompatible environments forces compaction rather than merely
favouring it. The size of the genetically encoded residue alphabet follows
from a purely local, mechanical substitution dynamic together with a
recurrence criterion for what counts as a storable category, with no step
at which the several hundred chemically available residues are compared
against one another by any evaluating process.

The final part of the derivation addresses a question the preceding results
raise but do not by themselves answer: given that a storage structure
exists, encoding a set of residues fixed by cross-lineage recurrence, what
determines when that structure is consulted. The answer given here, that a
local partition imbalance couples to a stored structure by a purely
structural and electrostatic criterion, blind to whatever product results,
is offered specifically because the alternative, in which the coupling is
mediated by an anticipated correspondence between stored structure and
product, presupposes exactly the correspondence the rest of the derivation
is constructed to explain rather than assume. A mechanism of reading that
requires foreknowledge of what it will produce cannot be the mechanism that
brought the correspondence between sequence and product into existence in
the first place, and the framework developed here is constructed throughout
to avoid this circularity: every result obtained is obtained from local,
mechanical, and retrospectively evaluated events, never from a step in
which some element of the system already knows what a later step will
require.

\section{Scope and Open Questions}
\label{sec:limitations}

The derivation of the encoded alphabet's size in \cref{part:alphabet}
identifies the mechanism that fixes the count and identifies the count as a
recurrence statistic, but it does not compute the numerical value
twenty-two from the combinatorics of any specific set of prebiotic
chemical pathways; doing so would require an inventory of the
self-propagating structures accessible under early conditions that is
beyond the scope of the present derivation and is identified here as
outstanding work. Similarly, while \cref{part:reading} establishes the
structural criterion by which a stored sequence is coupled to a local
imbalance, it does not derive the specific molecular apparatus, such as the
proteins that in extant biology mediate this coupling, from the partition
postulate alone; the claim advanced is that any such apparatus, whatever its
detailed chemistry, must satisfy the blindness condition of
\cref{def:blindcoupling} on pain of the circularity identified in
\cref{thm:blind-required}, not that the apparatus itself has been derived
in full molecular detail. The uniform floor of \cref{thm:uniform-floor}
remains, as flagged throughout, conditional on the bounded-refinement
hypothesis of \cref{def:res-bounded}, and every result inheriting that floor
inherits the condition with it.

\section{Conclusion}
\label{sec:conclusion}

A protein is conventionally introduced as a sequence of residues drawn from
a fixed alphabet of twenty amino acids, or twenty-two once the two
exceptional cases are included, folded into a compact structure by forces
that are catalogued but whose necessity is not established. We have shown
instead that the two-part architecture of the peptide backbone is forced by
the requirement that a bounded polymer present a spatially addressable
pattern of partition boundaries to its surroundings while remaining one
connected object; that folding is forced, not merely favoured, by the
impossibility of an extended chain presenting more than one partition
environment at once; and that the size of the encoded alphabet is fixed,
without any selecting or comparing agency, by how many residues happened to
recur as favourable local substitutions across independent self-propagating
lineages before a storage mechanism, itself explicable by the same
charge-emergence principle that governs every other partition boundary
considered in this paper, stabilised around them. Nothing in this account
requires that a genome, or any structure preceding it, possess knowledge of
the proteins it will come to encode; the correspondence between sequence and
product is, on this account, not a precondition of the mechanism but its
accumulated and entirely retrospective output.

\end{document}
